\section{Experimento: ¿Miel, azúcar o café?}

\subsection{Introducción}

En distintos sistemas de cultivo artesanal y huertos caseros suele mencionarse que algunas sustancias comunes pueden favorecer el crecimiento vegetal. Entre las más utilizadas aparecen la miel, el azúcar y el café soluble.

Frecuentemente se escuchan comentarios como:

\begin{quote}
“La miel contiene nutrientes y podría favorecer el crecimiento.”
\end{quote}

o bien:

\begin{quote}
“El café ayuda a estimular las plantas.”
\end{quote}

Sin embargo, muchas de estas afirmaciones provienen únicamente de observaciones empíricas y no de experimentos controlados.

Por esta razón se desarrolló un experimento utilizando semillas de lenteja con el objetivo de evaluar si distintos tratamientos generan diferencias observables en la germinación y desarrollo inicial de las plantas.

Además de estudiar el crecimiento vegetal, esta práctica permite introducir conceptos fundamentales del Diseño y Análisis de Experimentos, tales como:

\begin{itemize}
    \item unidad experimental,
    \item factor,
    \item niveles,
    \item réplicas,
    \item variabilidad experimental,
    \item análisis de medias,
    \item y ANOVA.
\end{itemize}

\subsection{Pregunta de investigación}

\begin{center}
\textit{¿Existen diferencias en la germinación y crecimiento de lentejas al aplicar miel, azúcar o café soluble?}
\end{center}

\subsection{Hipótesis}

\subsubsection{Hipótesis nula}

\[
H_0:\mu_B=\mu_M=\mu_A=\mu_C
\]

No existen diferencias significativas entre las medias de los tratamientos.

\subsubsection{Hipótesis alternativa}

\[
H_1:\text{Al menos una media es diferente}
\]

\subsection{Objetivo general}

Evaluar el efecto de distintos tratamientos sobre la germinación y crecimiento inicial de semillas de lenteja mediante un diseño experimental completamente aleatorizado y análisis ANOVA.

\section{Materiales y métodos}

\subsection{Materiales}

\begin{itemize}
    \item Agua
    \item Grenetina
    \item Azúcar
    \item Café soluble
    \item Miel
    \item Agua oxigenada comercial
    \item Antihongos comercial
    \item Semillas de lenteja
    \item Pipetas
    \item Báscula
    \item Recipientes plásticos transparentes
    \item Regla milimétrica
\end{itemize}

\subsection{Construcción del medio semisólido}

Para construir el medio de cultivo se utilizaron las siguientes cantidades:

\begin{table}[H]
\centering
\begin{tabular}{lc}
\toprule
Componente & Cantidad \\
\midrule
Agua & 2000 g \\
Grenetina & 100 g \\
Azúcar & 100 g \\
Agua oxigenada & 30 g \\
Antihongos & 0.5 g \\
\bottomrule
\end{tabular}
\caption{Composición del medio base}
\end{table}

\subsection{Preparación del medio}

\begin{enumerate}
    \item Calentar el agua.
    \item Agregar grenetina y azúcar.
    \item Llevar a ebullición.
    \item Mantener el hervor durante cinco minutos.
    \item Retirar del fuego.
    \item Esperar ligera disminución de temperatura.
    \item Agregar agua oxigenada y antihongos.
    \item Mezclar cuidadosamente.
    \item Colocar aproximadamente 30 g de medio por recipiente.
\end{enumerate}

\section{Tratamientos experimentales}

Cada solución fue preparada utilizando:

\[
30\text{ g de agua} + 30\text{ g de soluto}
\]

\begin{table}[H]
\centering
\begin{tabular}{ccc}
\toprule
Color & Tratamiento & Código \\
\midrule
Amarillo & Blanco & B \\
Morado & Miel & M \\
Azul & Azúcar & A \\
Verde & Café soluble & C \\
\bottomrule
\end{tabular}
\caption{Tratamientos experimentales}
\end{table}

\subsection{Aplicación experimental}

\begin{itemize}
    \item Se utilizaron 6 semillas por recipiente.
    \item Cada tratamiento fue realizado por triplicado.
    \item Las semillas fueron previamente hidratadas.
    \item Se aplicaron 3 gotas del tratamiento correspondiente.
    \item Los recipientes permanecieron cerrados.
    \item Las observaciones fueron realizadas diariamente.
\end{itemize}

\section{Diseño experimental}

\begin{itemize}
    \item Factor: tratamiento aplicado
    \item Niveles: 4
    \item Réplicas: 3
    \item Diseño: completamente aleatorizado
\end{itemize}

El número total de unidades experimentales fue:

\[
4\times3=12
\]

\subsection{Unidad experimental}

La unidad experimental corresponde al recipiente y no a cada semilla individual.

Esto es importante porque todas las semillas dentro de un mismo recipiente comparten:

\begin{itemize}
    \item humedad,
    \item temperatura,
    \item iluminación,
    \item y concentración del tratamiento.
\end{itemize}

\section{Variables evaluadas}

Durante el experimento se registraron:

\begin{itemize}
    \item Número de semillas germinadas
    \item Longitud de raíz
    \item Altura del tallo
    \item Presencia de hojas
    \item Presencia de hongos
    \item Condensación
\end{itemize}

\section{Estrategia temporal de análisis}

Aunque el experimento fue monitoreado diariamente, únicamente se reportarán medias correspondientes a:

\begin{itemize}
    \item Día 2
    \item Día 4
    \item Día 20
\end{itemize}

Estos tiempos representan distintas etapas del desarrollo experimental:

\begin{itemize}
    \item Día 2: germinación inicial,
    \item Día 4: establecimiento temprano,
    \item Día 20: desarrollo consolidado.
\end{itemize}

El análisis estadístico inferencial mediante ANOVA fue realizado exclusivamente utilizando los datos correspondientes al día 20, debido a que este punto temporal representa el estado final del sistema experimental bajo condiciones homogéneas de exposición.

\section{Tabla general de registro}

\begin{table}[H]
\centering
\resizebox{\linewidth}{!}{
\begin{tabular}{lcccccc}
\toprule
Tratamiento & Réplica & Día 2 & Día 4 & Día 20 & Hongos & Observaciones \\
\midrule
Blanco & 1 & & & & & \\
Blanco & 2 & & & & & \\
Blanco & 3 & & & & & \\
\midrule
Miel & 1 & & & & & \\
Miel & 2 & & & & & \\
Miel & 3 & & & & & \\
\midrule
Azúcar & 1 & & & & & \\
Azúcar & 2 & & & & & \\
Azúcar & 3 & & & & & \\
\midrule
Café & 1 & & & & & \\
Café & 2 & & & & & \\
Café & 3 & & & & & \\
\bottomrule
\end{tabular}
}
\caption{Registro experimental resumido}
\end{table}

\section{Cálculo de medias experimentales}

A partir de las mediciones registradas en las fotografías correspondientes al día 20 del experimento, se calcularon las medias de crecimiento para cada tratamiento.

En el tratamiento con miel se registraron las siguientes longitudes de raíz:

\[
4.2,\ 4.5,\ 5.1 \text{ cm}
\]

La media experimental fue calculada mediante:

\[
\bar{x}=\frac{\sum_{i=1}^{n}x_i}{n}
\]

Sustituyendo los valores experimentales:

\[
\bar{x}=\frac{4.2+4.5+5.1}{3}
\]

\[
\bar{x}=\frac{13.8}{3}
\]

\[
\bar{x}=4.6\text{ cm}
\]

\section{Cálculo de varianza}

La dispersión experimental fue estimada mediante la varianza muestral:

\[
s^2=\frac{\sum_{i=1}^{n}(x_i-\bar{x})^2}{n-1}
\]

Sustituyendo los valores observados:

\[
s^2=\frac{
(4.2-4.6)^2+
(4.5-4.6)^2+
(5.1-4.6)^2
}{2}
\]

\[
s^2=
\frac{
0.16+0.01+0.25
}{2}
\]

\[
s^2=
\frac{0.42}{2}
\]

\[
s^2=0.21
\]

\section{Análisis ANOVA}

El análisis ANOVA permite determinar si las diferencias observadas entre tratamientos son suficientemente grandes como para pensar que existe un efecto real asociado al tratamiento aplicado.

\subsection{Estadístico F}

\[
F=\frac{CM_{Tratamientos}}{CM_{Error}}
\]

donde:

\begin{itemize}
    \item $CM_{Tratamientos}$ corresponde al cuadrado medio entre tratamientos,
    \item $CM_{Error}$ corresponde al cuadrado medio del error experimental.
\end{itemize}

Si:

\[
p<0.05
\]

entonces se rechaza la hipótesis nula.

\section{Tabla ANOVA}

\begin{table}[H]
\centering
\begin{tabular}{lcccc}
\toprule
Fuente & SC & gl & CM & F \\
\midrule
Tratamientos & & & & \\
Error & & & & \\
Total & & & & \\
\bottomrule
\end{tabular}
\caption{Tabla ANOVA}
\end{table}

\section{Interpretación biológica}

Si el ANOVA no presenta diferencias significativas, entonces las variaciones observadas entre tratamientos podrían atribuirse al azar experimental.

Por el contrario, si el análisis presenta diferencias significativas, entonces al menos uno de los tratamientos estaría afectando el crecimiento vegetal de manera distinta.

En ese caso podrían realizarse pruebas post-hoc, como Tukey, para identificar específicamente qué tratamientos difieren entre sí.

\section{Conclusiones}

El experimento permitió integrar observación biológica, medición experimental y análisis estadístico dentro de un mismo sistema de estudio.

La utilización de tratamientos como miel, azúcar y café soluble permitió construir un escenario experimental adecuado para introducir conceptos fundamentales del Diseño y Análisis de Experimentos.

El uso de réplicas permitió estimar variabilidad experimental, mientras que el análisis ANOVA proporcionó una herramienta formal para comparar tratamientos bajo condiciones homogéneas.

La decisión de utilizar únicamente las mediciones correspondientes al día 20 para el análisis inferencial permitió homogenizar temporalmente las comparaciones y reducir problemas asociados a mediciones repetidas dependientes.

Finalmente, el experimento muestra cómo observaciones aparentemente simples pueden transformarse en un análisis científico estructurado mediante:

\begin{itemize}
    \item control experimental,
    \item medición,
    \item estadística,
    \item y análisis cuantitativo.
\end{itemize}